\section{Artifacts and Signal Quality}

Wearable EEG during game interaction is vulnerable to eye movement, blinks, facial and jaw electromyography, head motion, band displacement, poor contact, packet loss, and Bluetooth buffering. These problems are part of the measurement question and must be reported rather than hidden by selective retention of clean events.

The study freezes rules for minimum usable channels/features, maximum packet loss, amplitude clipping, movement contamination, minimum clean duration, and window/channel rejection before inferential analysis. Accelerometer and gyroscope streams should be retained as artifact/context signals when the authorized Athena research interface exposes them.

Low-gamma power is treated as secondary because facial and cranial muscle activity can overlap the relevant frequency range. No cortical source-localization claim is made from the sparse wearable montage.

Signal-quality exclusions must be independent of whether the behavioral outcome is favorable. The paper reports the full candidate-event denominator, quality-failure reasons, and the extent to which missing EEG is associated with movement, challenge intensity, or participant characteristics.
